% Authorship: Claude-drafted, author-revised
\section{Berman's Principled Positivism}

This Part reconstructs Berman's theory with the precision that
computational implementation demands.\footnote{This reconstruction draws
primarily from Berman, \textit{Practices}, \textit{supra} note 7.
Berman develops the same principles-based account of legal content in two
further works on which I also draw: Mitchell N.\ Berman, \textit{Our
Principled Constitution}, 166 \textsc{U.\ Pa.\ L.\ Rev.}\ 1325 (2018)
(arguing that the content of American constitutional law is fixed by a
plurality of legal principles rather than by rules alone), and Mitchell
N.\ Berman, \textit{Religious Liberty and the Constitution: Of Rules and
Principles, Fixity and Change}, 26 \textsc{U.\ Pa.\ J.\ Const.\ L.}\ 851
(2024) (showing how a plurality of principles can yield singular
constitutional answers despite their plurality, and addressing the
resolution of principle conflict beginning at 889).} Throughout, the goal is to operationalize Berman's philosophical arguments:
to specify the model's claims, the parameters it requires, and the
questions it leaves unresolved. For readers already familiar with the
theory, the Section to attend to is likely the last one. It identifies the
three points where the theory leaves its content open and singles out the
one that motivates the computational analysis.

\subsection{The Core Model}

Berman's aggregation model comprises three conceptually distinct
elements: \textit{principles}, \textit{weights}, and
\textit{activations}.\footnote{Berman, \textit{Practices}, \textit{supra} note 7, at 328--36.} Principles are legally
fundamental norms that contribute to determining outcomes in particular
cases. In two respects that matter for the aggregation model, they differ from
rules.\footnote{Berman, \textit{Practices}, \textit{supra} note 7, at 299--303. The principle/rule distinction
has a long pedigree. See Ronald Dworkin, \textit{The Model of Rules}, 35
\textsc{U.\ Chi.\ L.\ Rev.}\ 14, 25--29 (1967). Berman's account of how
principles \textit{operate}, however, differs substantially from
Dworkin's: for Berman, principles contribute weighted force through
aggregation; for Dworkin, they provide reasons requiring interpretive
judgment.} First, principles admit of degree. Where state action directly degrades
personhood, ``respect human dignity'' applies strongly; where the
connection is attenuated, weakly; where dignity is simply not implicated,
not at all. By contrast, rules are categorical: their conditions either
obtain or they do not. Second, principles can conflict without
contradiction. When textual plain meaning opposes an outcome and human
dignity favors it, both remain valid legal norms; they pull in opposite
directions rather than generating logical inconsistency. Without that second feature, aggregation would not be possible.

In U.S.\ constitutional and statutory law, examples include respect for
textual plain meaning, honor for precedent, fidelity to original meaning,
protection of human dignity, promotion of equal protection, protection
of individual liberty, preservation of federalism, avoidance of absurd
results, and promotion of general welfare.\footnote{These examples draw
from Berman, \textit{Practices}, \textit{supra} note 7, at 330--36 and are
illustrative rather than
exhaustive. The computational model uses nine principles of these kinds,
though not this exact list; Part III specifies the set.}

Principles are \textit{constituted} by practice. On Berman's account, a norm counts as a legal principle not because it is
morally true but because legal officials \textit{take it up}, systematically
invoking it in their reasoning.\footnote{Berman, \textit{Practices}, \textit{supra} note 7, at 329 (``Principles
are constituted by the taking-up behaviors of the system's
subscribers.''). This is what makes Berman a positivist: law's content
depends on social facts about official practice, not on moral facts about
what justice requires.} In that sense principled positivism remains positivist, despite trafficking
in norms like dignity, equality, and liberty that look like moral values.
On both counts, its source in practice and its independence from moral
truth, the account differs from Dworkin's, where
principles are objects of moral interpretation,\footnote{Dworkin, \textsc{Law's Empire},
\textit{supra} note 4, at 225--75.} and from natural law theories, where
they are moral truths accessible to reason.\footnote{See generally Mark C.\
Murphy, \textsc{Natural Law in Jurisprudence and Politics} (2006).}

The set of principles is not fixed. As judicial practice changes, so does
the set: principles can enter the legal system as new norms gain traction
and can exit as old ones lose it.\footnote{Exit is more
complicated than mere disuse. A principle is embedded by the holdings that
rely on it, so it can retain vitality even when it is not currently being
invoked; and a principle may be affirmatively \textit{disavowed}, which
Berman distinguishes from its merely falling into desuetude. See infra
Section~II.D (developing these points, which bear on how the dynamics of
practice are modeled).} Whether a norm qualifies as a legal principle in
the first place depends on a \textit{threshold of uptake} (not the
determinacy threshold introduced below): enough officials must
invoke the norm for it to cross from a merely advocated position to a genuine
legal principle.\footnote{Berman, \textit{Practices}, \textit{supra} note 7, at 329.} That uptake threshold, too, is presumably practice-grounded.

Weights represent each principle's context-invariant importance. Each
principle has a weight $w_i \in [0,1]$. On such a scale, textual plain meaning might carry weight 0.85, reflecting
its prominence across interpretive methodologies;
even non-textualists acknowledge that plain meaning
matters.\footnote{See William N.\ Eskridge Jr.\ \& Philip P.\ Frickey,
\textit{Statutory Interpretation as Practical Reasoning}, 42 \textsc{Stan.\
L.\ Rev.}\ 321, 353--57 (1990).} By the same measure, federalism might carry weight 0.65, reflecting its
less central but still significant role.\footnote{Berman, \textit{Practices}, \textit{supra} note 7, at 328--30. These illustrative values are not empirical claims about actual
U.S.\ law. The computational analysis tests performance across diverse
weight distributions precisely because the ``true'' values are
unknown.} On Berman's account the weights are determined by practice: the more
centrally a principle figures in legal reasoning across diverse cases,
methodologies, and time periods, the higher its weight.\footnote{Berman, \textit{Practices}, \textit{supra} note 7, at 329.}

Activations represent how strongly a principle applies in a particular
case. Formally, $a_i \in [-1, +1]$, where the magnitude indicates
strength and the sign indicates direction: positive favors the proposed
outcome, negative opposes it.\footnote{Berman, \textit{Practices}, \textit{supra} note 7, at 330--32, 339--41.} When
statutory text clearly supports a particular reading, the textual
principle is strongly activated ($a \approx +0.9$). When the text is
ambiguous, activation is weak ($a \approx +0.3$). When the statute does
not address the question at all, activation is zero. Crucially, the sign captures the idea that a principle can favor
\textit{either side} of a dispute.

The distinction between weight and activation is central to the model and
easily lost. A weight is constant across cases and an activation varies with them:
the relative importance of precedent does not change from case to case, even
though precedent's applicability does.\footnote{A judge who says ``precedent
matters more in this case'' is, on Berman's account, saying that precedent is
more strongly \textit{activated}, not that its \textit{weight} has changed.}

Activation is where most disagreement tends to occur in hard
cases.\footnote{Berman, \textit{Practices}, \textit{supra} note 7, at 339 (``Perhaps the justices disagree about these
principles' relative weights. More conspicuously and consequentially,
however, they disagree about the extent to which each principle was
activated.'').} When Justices divide in a statutory interpretation case,
they typically agree about \textit{which} principles matter and roughly
\textit{how important} each one is, and divide instead over \textit{how
strongly} each principle applies to the facts, though the agreement on
weight is a tendency, not a rule.\footnote{Berman does not claim the
disagreement is ever \textit{solely} about activation. In \textit{TVA v.\
Hill}, discussed below, he reads the Justices as having differed slightly
about the \textit{weight} of the avoid-absurdity principle as well, not
only its activation. Berman, \textit{Practices}, \textit{supra} note 7, at 338--42 \& figs.\ 8--9.} Here the Article's contribution comes into view. If officials largely
converge on the weights and divide over activation, then the pressing
question for the theory is likely where the weights come from at all: the
question Parts V and VI take up.

Given these three elements, the model computes net legal force through
Simple Additive Weighting: $F = \sum_{i=1}^{n} w_i \cdot a_i$. A case is
\textit{determinate} when the principles, taken together, favor one
side, that is, when the balance is not in equipoise. To that one may add a
\textit{determinacy threshold} $\tau$, counting a case as determinate only
when $|F|$ exceeds $\tau$ and as a genuine legal gap when $|F| \leq \tau$.
But as the last Section of this Part explains and Part VIII confirms, that
threshold is a modeling choice rather than a commitment of the theory, and the faithful
baseline is simply the sign of net force ($\tau = 0$).\footnote{Berman, \textit{Practices}, \textit{supra} note 7, at 316--18, 328--31. Simple Additive Weighting is a standard technique in
multi-criteria decision analysis. See Ralph L.\ Keeney \& Howard Raiffa,
\textsc{Decisions with Multiple Objectives: Preferences and Value
Tradeoffs} 75--98 (1976).}

\subsection{A Worked Example: \textit{TVA v.\ Hill}}

To see how the model operates on a real dispute, consider \textit{TVA v.\
Hill}.\footnote{437 U.S.\ 153 (1978).} There the Supreme Court divided
six-to-three over whether the Endangered Species Act required halting
construction of the nearly completed Tellico Dam to protect the snail
darter. Berman discusses this case extensively as a paradigm of principled
disagreement,\footnote{Berman, \textit{Practices}, \textit{supra} note 7, at
338--42.} and it repays careful analysis, since it illustrates every component
of the model in action, perhaps better than any other case.

Section 7 of the ESA directs all federal agencies to ``insure that actions
authorized, funded, or carried out by them do not jeopardize the continued
existence'' of endangered species.\footnote{16 U.S.C.\ \S~1536 (1976 ed.);
\textit{quoted in} 437 U.S.\ at 173.} Chief Justice Burger's majority held
that it flatly prohibited completing the dam. Justice
Powell dissented, arguing that the statute was not so clear as applied to
a project substantially complete when the species was listed.

On Berman's reading, Burger and Powell agree about \textit{which}
principles matter and differ only slightly about \textit{how important}
they are. Both
recognize that textual plain meaning, congressional intent, and avoiding
absurd results bear on the question, and neither contests their relative
ranking. The disagreement is about activation: how strongly each
principle applies to \textit{this particular case}.

Burger reads the statutory text as maximally plain. ``One would be hard
pressed,'' he writes, ``to find a statutory provision whose terms were any
plainer than those in \S~7.''\footnote{437 U.S.\ at 173.} Finding congressional intent unqualified, he treats the costly result as
Congress's considered policy choice rather than an absurdity. Powell
reads the same materials very differently. He describes the critical
statutory term ``actions'' as ``a textbook example of fuzzy
language.''\footnote{\textit{Id.}\ at 202 (Powell, J., dissenting).} He
finds legislative intent silent on nearly completed projects. And he
regards halting the dam as manifestly absurd.

Table~\ref{tab:tvah} formalizes this disagreement, holding the weights fixed
and varying only the activations. A positive activation favors prohibition,
a negative one favors completion.

\begin{table}[ht]
\centering
\caption{Activation Disagreement in \textit{TVA v.\ Hill}}
\label{tab:tvah}
\begin{tabular}{lccc}
\toprule
\textbf{Principle} & \textbf{Weight} & \textbf{Burger} & \textbf{Powell} \\
\midrule
Textual plain meaning & 0.85 & $+0.9$ & $+0.3$ \\
Congressional intent & 0.70 & $+0.8$ & $0.0$ \\
Avoid absurdity & 0.60 & $-0.2$ & $-0.9$ \\
\midrule
\textbf{Net Force} & & $\mathbf{+1.205}$ & $\mathbf{-0.285}$ \\
\bottomrule
\end{tabular}
\end{table}

Burger's analysis yields strong positive net force ($+1.205$): the law,
on his reading, clearly favors prohibition. His textual figure alone
contributes $0.85 \times 0.9$, or $0.765$. Powell's analysis yields $-0.285$: a weak textual contribution, no intentionalist
contribution at all, and a strong anti-absurdity force that carries the
balance the other way. On his reading the statute does not prohibit
completion, the conclusion he reached.\footnote{The zero for
congressional intent is Berman's. He renders Powell's view of legislative
intent in Figure 9 as a double-headed arrow, pointing toward prohibition
and away from it at once, which is his diagram of the conclusion that
``there was no actual congressional intention relevant to this dispute, no
intention either that completion of the Tellico Dam project would be
illegal or that it would not be.'' Berman, \textit{Practices},
\textit{supra} note 7, at 340; \textit{see id.}\ at 342 figs.\ 8--9. In a
signed scalar model a principle pulling equally both ways contributes zero,
which is the value used here. That is not the same state as a principle
nobody has engaged, and the model cannot distinguish them: both enter the
sum as $0$. The limitation is worth naming, since Berman's own diagram
draws a distinction the scalar representation cannot carry.}

Predominantly, though not exclusively, the disagreement is about activation.
As noted above, Berman reads the Justices as having also
differed slightly about the weight of the avoid-absurdity principle.\footnote{Berman, \textit{Practices}, \textit{supra} note 7, at 338--42 \& figs.\ 8--9.} Even so, the case illustrates Berman's claim that activation, not weight,
drives most judicial disagreement.

The example also shows what the determinacy threshold does, and does not,
control. Powell's near-equipoise reading sits close to zero: on the
faithful sign-of-force baseline the negative sign of his net force favors his
anti-prohibition conclusion, while a modest positive threshold would
instead make his case a genuine gap. Here, then, the threshold would govern how much of the law is reported as
determinate, an effect Part VIII examines. It leaves untouched the
net forces themselves and the substantive properties that flow from them.

\subsection{Three Claims}

Berman advances three substantive claims about the model's performance.
Each responds to a specific objection to principled positivism, and each
makes an empirically testable prediction about how the model should
behave.

The first claim is that rough weights suffice. Stephen Perry pressed
perhaps the most natural objection to any aggregation model: if the
weights come from the messy, contested practice of diverse
judges, how can they be precise enough for aggregation to produce
determinate results?\footnote{Perry, \textit{supra} note 9, at 794.}
Berman's response invokes what might be called the \textit{sheer numbers
effect}: when multiple principles point in the same direction, their combined force
should swamp small errors in individual
weights.\footnote{Berman, \textit{Practices}, \textit{supra} note 7, at 328--30.} If five principles favor an outcome and three oppose it, getting the
weights exactly right matters less, arguably, than getting the rough
balance right. As an empirical claim about the model's tolerance for weight uncertainty,
this one is directly testable, and Part IV tests it by systematically
perturbing weights and measuring how much imprecision the model can absorb.

The second claim is that conflicting principles coexist stably. A Hartian
validation rule, one that asks of each standard only whether it is valid by
the system's criteria and which way it points, faces a problem when
conflicting standards are both valid.\footnote{See Shapiro,
\textsc{Legality}, \textit{supra} note 6, at 270--75.} If textualism and
dignitarian reasoning both satisfy the rule's criteria for legal validity
and point in opposite directions in a given case, the validation rule
appears to generate contradictory law: the case is both legally required
and legally prohibited.\footnote{For an influential analysis of this
problem, see Mark Greenberg, \textit{How Facts Make Law}, 10 \textsc{Legal
Theory} 157, 185--89 (2004).} Berman's model dissolves the problem.
Conflicting principles do not generate contradictions; they generate
opposing forces that sum algebraically. Both principles remain valid and operative, and simply pull in opposite
directions. In Part~IV that contrast sharpens into an exact one. A rule
that asks only whether each standard is valid contradicts itself on a
genuine conflict; the weighted sum never can. As a description of how
judges actually reason, the outweighing structure is arguably the more
recognizable of the two. In \textit{Obergefell v.\ Hodges}, Justice Kennedy
did not claim that originalist arguments against same-sex marriage were
\textit{invalid}; he claimed they were \textit{outweighed}.\footnote{576
U.S.\ 644, 666--76 (2015).} Berman's model formalizes this outweighing
structure, and Part~IV tests it: cases with balanced conflict should still
produce determinate outcomes at meaningful rates.

The third claim is that the model delivers more law than Hart. Of the three, this one most directly addresses the Hart-Dworkin impasse. Berman
argues that aggregation solves Dworkin's ``too little law'' problem:
principles that individually lack consensus support can aggregate to
determinate outcomes when their combined force is sufficient.\footnote{Berman, \textit{Practices}, \textit{supra} note 7, at 316--18.} Several principles, each with partial
support, can add up to determinate law where Hart finds only a gap. On this the prediction is straightforward. Berman's model should produce
higher determinacy rates than a Hartian baseline.

The baseline here is not the validation rule described above but a more
charitable construction, a \textit{hierarchical Hart}. It sorts the
principles into ranked tiers by kind, textual above structural above moral,
and decides a case at the highest tier that contains an applicable principle,
provided enough of that tier's principles point the same way. Where they do
not, it leaves the case unsettled and never reaches the tiers below, and so
gives the weighted sum something to beat. Part~IV measures how large the
margin over it turns out to be.

These three claims bear on three of the four requirements this Article tests
the model against, and Part III states all four precisely. The
coexistence claim bears on the first, that the weights be
\textit{well-defined}, that is, that a plurality of principles carry
genuine weight at all. The rough-weights claim bears on the second, that they
be \textit{noise-tolerant}. The more-law claim bears on the fourth, that the
model deliver \textit{more law than Hart}. The third requirement comes not
from any of the three claims but from the theory's rejection of lexical
priority: that the model be \textit{magnitude-weighing}, responding to how
strongly each principle applies rather than merely to which principle ranks
highest. Rejecting lexical priority is not by itself enough to satisfy it,
since a rule that counts principles rejects lexical priority too and still
reads no magnitudes. That one occupies Part VII. Two further properties that earlier work has
associated with the theory (that principles travel in correlated groups,
and that their weights come graded rather than equal) are worth measuring
but are \textit{not} requirements. Berman's account, unlike Dworkin's,
permits principles to carry equal weight.\footnote{Berman, \textit{Practices}, \textit{supra} note 7, at 323--24 (the
account does not require that principles have varied weights). Equal
weights do not collapse the model into majority voting, a proposition Part
IV establishes formally.}

\subsection{What the Theory Leaves Open}

Berman's model is precise for a philosophical theory. It specifies the
aggregation equation, identifies the relevant parameters, and makes
predictions about the model's performance. Of the four questions the
Introduction left open, the first is Berman's own and he answers it: rough
weights suffice. How rough is a quantity to be measured rather than a gap to
be filled, and Part III accordingly states it as a requirement. Yet
translating the model into executable code, which demands full
specification, reveals three further places where the theory does not fully
determine its own content. They are not equally consequential. Two of the three look open but would seem to settle largely on inspection.
The third does not, and it is where this Article's central finding lies.

\subsubsection{The Determinacy Threshold}

Where does $\tau$ come from? The aggregation equation specifies how to compute
net legal force; the determinacy condition specifies when that force is
\textit{sufficient}, and that second specification looks as though it needs a
threshold the theory never grounds. But Berman does not, in fact, commit
to one. He writes that an aggregative system ``could incorporate
thresholds or eschew them,''\footnote{Berman, \textit{Practices}, \textit{supra} note 7, at 331.} and in correspondence he
has confirmed that he claims no more than that some threshold is
\textit{possible}, not that any is \textit{required}.\footnote{E-mail from
Mitchell N.\ Berman to author (Mar.\ 2026) (on file with author). Berman
observes that, even with no threshold, not every dispute would have a
determinate answer, since the contours and weights of some principles, and
their context-specific activation, may themselves be indeterminate, but
that where the principles and their activation \textit{are} determinate,
there is law ``whenever the balance of principles is not precisely in
equipoise,'' with some actions ``barely legal'' or ``barely illegal.''}
The faithful baseline, then, is the bare \textit{sign} of net force: there
is law wherever the balance favors one side. A nonzero threshold is not a commitment of the theory but a
\textit{modeling choice}. On the evidence of Part VIII, it governs how much
of the law comes out determinate without disturbing the model's substantive
properties.

\subsubsection{The Aggregation Function}

Berman's model sums weighted activations. Is that the right functional
form? Because the theory already takes a side on the only distinction that
matters here, the question has a narrower compass than it first appears.
Berman's principles combine by \textit{weight and magnitude}: a strongly
activated principle can be outweighed by several weaker ones pulling the
other way. They do not combine by \textit{fixed priority}, under which a
higher-ranked principle would prevail over any lower-ranked one regardless
of magnitude. The first mode is \textit{nonlexical}; the second is
\textit{lexical}, and Berman is explicit that his account rejects
it.\footnote{Berman, \textit{Practices}, \textit{supra} note 7, at
311--13, 333 (distinguishing genuine weighing from
lexical priority and placing principled positivism on the weighing side of
that line). Lexical priority is associated with Joseph Raz's
exclusionary-reasons view. See Joseph Raz, \textsc{Practical Reason and
Norms} 35--48 (2d ed.\ 1999).} Lexicographic and cascade orderings let the
highest-ranked applicable principle decide and consult nothing below it. On
that account they are not rival readings of principled positivism at all.
They are the foil it is defined against, useful here only as a contrast.

Ruling out lexical priority still leaves a choice, because the nonlexical
rules divide again. Some read how strongly each principle applies; others
read only which way each one points and count. The live question is
therefore twofold. Within the rules that read strength, does anything turn
on the particular form, simple addition against, say, a combination in
which a principle's contribution grows faster than linearly? Part VII finds
that nothing does. Across the line separating those rules from the ones
that merely count, a great deal turns on the choice: the rules that read
strength and the rules that only count reach different verdicts on about one
case in seven. So what does the work is not
nonlexicality, which the counting rules share, but the reading of
magnitude, which they do not.\footnote{Berman draws the additive form
from multi-criteria decision analysis, where simple additive weighting is
the default. Berman, \textit{Practices}, \textit{supra} note 7, at 330; see
Keeney \& Raiffa, \textit{supra} note 32, at 75--98.}

\subsubsection{The Grounding of the Weights}

This is the deep gap, and it is the question the computational analysis is
built around. Berman says the weights are grounded in practice: the more centrally a
principle figures in legal reasoning, the higher its weight. This Article
compresses that phrase to \textit{practice-grounded}, a term that is its own
and not Berman's.\footnote{Berman, \textit{Practices}, \textit{supra} note 7, at 329.} But \textit{how} does a pattern of practice fix
a number between zero and one? Here a distinction Berman has drawn is
essential. The grounding \textit{relation}, the source of a principle's
weight, is the taking-up and embedding of the principle by legal officials.\footnote{Berman, \textit{Practices}, \textit{supra} note 7, at 329--30; see also
id.\ at 321--22. Berman has emphasized, in
correspondence and in his book in progress, that a principle is also
embedded by the holdings that continue to rely on it, so that its vitality
does not depend on current invocation alone, and that a principle may be
\textit{disavowed}, which is distinct from its merely falling into
desuetude. See Berman, \textsc{To Say What the Law Is}, \textit{supra}
note 7 (manuscript ch.\ 6).} That much is the theory's answer, gestural though it is. The theory does
\textit{not} address a distinct, prior question: what \textit{drives changes} in those takings-up over time, the
\textit{dynamics of changing grounds}.

It is those dynamics this Article models, and it models not one but four.
A principle's weight may shift with the sheer \textit{uptake} of the
principle in judicial reasoning; with its \textit{embedding} in holdings
still relied upon; with the \textit{strategic adaptation} of judges who
learn which arguments move their colleagues; and with the near-random
\textit{exogenous shock} of judicial appointment, as new judges arrive
with their own views of which principles matter and how
much.\footnote{The appointment dynamic responds to Berman's suggestion,
in correspondence, that cultural uptake and strategic adaptation are not
the only dynamics at work and that appointment is a third. Berman also
cautions against the language of biological ``evolution'' and ``fitness''
in this setting, since there are few hard constraints on which approaches
survive beyond their social acceptance, a caution this Article adopts by
speaking of the dynamics of changing grounds rather than of selection for
fitness.} Part III implements the four as the uptake, reliance, strategic, and
appointment mechanisms, in that order, so the embedding dynamic appears there
and thereafter under the name reliance, and the shock of appointment simply as
appointment. These four dynamics are not equally distinct. The weight structure
appointment produces tracks embedding's so closely that the two are in
effect one. Appointment is also the weakest of the four, clearing the
noise-tolerance requirement only marginally. On the evidence of Part VI,
then, the four dynamics realize three distinct weight structures rather than
four. Whether even those three weight structures
ground the weights in the same way, the theory does not say, and that is
the question the rest of the Article addresses.

\subsubsection{Why the Grounding Question Is the One That Matters}

One might expect these gaps to compound, so that a theory underdetermined
on its threshold, its functional form, \textit{and} its grounding could
not deliver determinate law at all. The computational results point the
other way. The threshold governs only how much of the law is reported as
determinate, not whether the model's substantive properties hold. The functional form does not matter among the rules that read magnitude, the
family the theory endorses; what matters is whether a rule reads magnitude at
all. And the formal apparatus, given practice-grounded weights, does
everything Berman asks of it. The slack that does matter lies in the
grounding. Because the dynamics confer different weights, they divide
the cases differently.\footnote{Part VI reports how often, and on which
kinds of cases.} For that reason the choice among them is not a modeling convenience but a
determinant of legal content. In just that sense principled positivism is
underspecified. Not in its formal structure, which works, but in its
dynamics of changing grounds, from which the structure borrows its content
and which the theory has yet to supply. Establishing that requires a
framework first: an implementation of the aggregation, four implementations
of the practice that might ground it, and the requirements both are
measured against. Part III builds it. Parts V and VI then make the point
concrete.
